% !TEX root = ../main.tex
\lecture{2}{2026-09-21}{Woche 2}

\section{Bilanz}

\begin{itemize}
    \item Vermoegenslage (Seiten)
    \subitem Verm. gegenstände <=> Verpflichtungen
    \item  Jede Org. hat eine Bilanz (von privat, staatlich bis NGOs)
    \item Bildet ab oder protokolliert Verschiebungen der Vermögenslage einer Organisation.
\end{itemize}

\subsection{Aktiv vs passiv Konten}
- Beide Seiten der Bilanz muessen Gleich sein

\begin{definition}
    [Aktiver Vermögensstand] bedeutet möglicher zukünftiger Nutzens\textbf{\textit{zu}}fluss. Aktuell in Gebrauch vom Unternehmen. Hoher aktiver Wert bedeutet gute zukünftige Lage \footnote{Hier sieht man den Grund für den Begriff "zukünftig". So kann der Aktienwert wahrscheinlich schneller ermittelt werden.}
\end{definition}
\begin{definition}
    [passiver Vermoegensgegenstand] bedeutet möglicher zukünftiger Nutzens\textbf{\textit{ab}}fluss. Daraus entspringen die Unternehmensaktiven aus FK oder EK. Hoher Passiver Wert bedeutet schlechtere zukünftige Lage
\end{definition}

\begin{itemize}
    \item Aktivum (Verm. gegenstand)
        \subitem Kasse (Weil bargeld zukünftigen Nutzenszufluss verursachen kann)
        \subitem Maschinen (Nutzenszufluss aus der Maschine indem es produziert)
        \subitem Vorräte    
        \subitem Forderungen aus Lieferungsleistungen (Debitoren, offene \textit{positive} Rechnungen)
    \item Passivum (Verpflichtungen, oder \textbf{FREMDKAPITAL})
        \subitem Hier ist die SOll Seite \textbf{Negativ} und die Haben Seite \textbf{Positiv} 
        \subitem Bankschulden 
        \subitem Hypothekarschuld
        \subitem Verbindlichkeiten an Lieferungen und Leistungen (Kreditoren)
    \end{itemize}


\subsection{Eröffnung der Aktiv- Passivkonten}
\textbf{3. Schritte des Vorgehens}
\begin{enumerate}
    \item  Ausgangslage => Eröffnungsbilanz per 1.1.2001
    \begin{itemize}
        \item Aktiv
        \begin{itemize}
            \item Konto Kasse
        \item \textbf{Debitoren} 
        Hier werden nur erwartete Beträge im Soll + aufgeschrieben, wenn rechnung bezahlt dann ins Haben (-)
        \item Waren
        \item Mobilien
        \end{itemize}
        \item Passiv
        \begin{itemize}
            \item Kreditoren
        \end{itemize} 
    \end{itemize}  
    \item Uebertrag der Bestände der  Eröffnungsbilanz in die Konten
        \subitem Konto fuehrung mit jeweil 1. Eintrag des ABs fuer jedes Konto aktiv und passiv
    \item Verbuchung der Buchungstatsachen
    \begin{enumerate}
        \item Welche Konten sind betroffen? (Aktiv, Passiv, dann Debit oder Kreditoren usw.)
            \subitem \textbf{Immer 2 Konten} weil sonst keine ausgleichung + und - ins Gesamt \footnote{Das heisst: Die Werte für den Tausch, die Finanzierung oder Definanzierung existierten bereits und kamen aus einer internen oder externen Quelle/Konto (Kasse, EK, FK, usw.).\textbf{ Kein neuer Wert oder Mehrwertschöpfung}, erst bei Erfolgsrechnung}
        \item Wo werden die Beträge eingetragen?
            \subitem Soll oder Haben
        \item Buchungssatz
        \begin{itemize}
            \item Name der Buchungstatsache 
            \item Konto, das im Soll beeinflusst wird
            \item Konto, das Haben beeinflusst wird 
            \item Der Betrag
            \item Bsp. Fahrzeugkonto /(an) Kassekonto => 40000

        \end{itemize}
        \item Kontrolle -> Jederzeit das $\Sigma Soll = \Sigma Haben$
    \end{enumerate}
\end{enumerate}

\subsection{Buchungsvorgänge innerhalb der Bilanz}

\textbf{Vier Grundtypen von Bilanztransaktionen}
\textbf{Format} NR. - Soll/Haben - Betrag
\begin{itemize}
    \item Aktiventausch (Nullsumme)
    \subitem In der Bilanz ist auf der Aktiv Seite eine Zunahme und Abnahme im gleichen Betrag zu sehen.
    \item Passiventausch (Nullsumme)
        \subitem In Bilanz nur Passivseite tangiert im gleichen Betrag ZU und Abnahme zweier Konten
        \subitem Buchungs satz (- bei passiv links) Darlehen / Eigenkapital 40000 
        \subitem Schulden werden durch Eigenkapital ersetzt 
        \footnote{ im Bsp. Darlehen wurden Schulden in Eigenkapital verwandelt, das heisst der Kreditor wurde zum Eigentümer der Kapitals im Betrag von 40000 (Oft in Finanzkrisen angewendet sodass die Bank im Unternehmen intervernieren kann}
        \footnote{Beide Tauschtransaktionen zeigen wie bestehende Werte im Unternehmen bewegt worden sind }
    \item Bilanzzunahme (Finanzierung)
        \subitem  Trasaktion aktiv zu passiv Konto 
        \subitem Vom aktiven Soll (+) ins passive Haben (+)
        \subitem Z.B kauf einer Maschine gegen REchnung (Machinen/Kreditoren)
    \item Bilanzabnahme (Definanzierung)
        \subitem Trasaktion von passiv zu aktiv Konto 
        \subitem Vom passiven Soll (-) ins aktive Haben (-)
\end{itemize}
\pagebreak
\subsection{Aufbau und Struktur der Bilanz } 
\textit{siehe Folien OR, Art. 959a}

\textbf{Aktiven}
\begin{itemize}
    \item Soll (Aktiven nach Liquiditätsprinzip)\footnote{Rotationsgeschwindikeit des Kapitals von Ware oder Wertpapier in Geld}
        \subitem Aktiv UV Konten 
        \subitem Summe der Umlaufvermögen (UV) < 1 Jahr
        \subitem Aktiven AV Konten
        \subitem Summe der Anlagevermögen (AV) > 1 Jahr
        \subitem Total summe
\end{itemize}
\textbf{Passiven}
\begin{itemize}
    \item Haben (Passiven nach Fälligkeit)
        \subitem summe kurzfristiges Fremdkapital Konten 
        \subitem summe Langfristiges Fremdkapital
        \subitem Eigenkapital 2 mal weil Eigenkapital Konto und Passivkategorie 
\end{itemize}



\subsection{ Konten: Passiv oder Aktiv?}
\textit{Siehe S. 35-39} \url{file:///C:/Users/kalel/Desktop/uzh/unterlagen/FIN_ACC/Finanz.%20Rechnungswesen%20Buch.pdf} 
\begin{itemize}
    \item \textbf{Aktiv}
    \begin{itemize}
        \item Kasse (aktiv, Umlaufverm.)
        \item Vorräte Aktiv, Umlaufverm
        \item Fahrzeug A, AV
        \item Debitoren (aktiv, Uv)
        \item Bankguthaben A, Uv
        
        \item Mobilien Aktiv, AV
        \item Immobilien Aktiv, Anlagevermoegen
    \end{itemize}
    \item \textbf{Passiv}
    \begin{itemize}
        \item Bankkonto-Kontokorrent (Schuld): Passiv, k. Fr
        \item Kreditoren, P krz FK 
        \item Darlehensschuld Passiv, Langfristiges FREMDKAPITAL
        \item Eigenkapital P, EK (wieso passiv?)\footnote{Weil Kapital den Eigentümerpersonen gehört und zunkünftig durch Eigentümer entzogen werden kann}
    \end{itemize}
\end{itemize}

\subsection{Fazit}
In der Bilanz ist die gleichheit von Aktiven und Pasiven Werte immer erhalten, also ist in der Bilanz keine Verbesserung der zukünftigen Vermögenslage möglich. zunkünftige Nutzenszuflusse werden immer durch Nutzenabflüsse ausgeglichen.
Einseitige Zunahmen von Aktiven oder Passiven nur in Erfolgsrechnung.
\[
    Eigenkapital = \textit{Aktivvermögen} - Fremdkapital
\]
oder
\[
   = \textit{Residualaktivvermögen nach Reduktion der Verbindlichkeiten an Dritten}
\]